\section{Sufficiency of the Conversion Rule Within the Model}\label{sec:sufficiency}

The frozen cohort shows the conversion rule is \emph{necessary} at
mechanism level: every measured policy that retires before
converting fails the acceptance contract. This section proves the
complementary direction---\emph{sufficiency}---for the frontier policy
under the engine's semantics, so the paper's rule is
necessity-plus-sufficiency within the model rather than an empirical
regularity. The proposition is scoped to the deterministic engine and
to streams with reliable retrieval (every fetch succeeds and returns a
current value); the unreliable-retrieval boundary is exactly where the
rule's guarantees stop, and Section~\ref{sec:retire-results} measures
what bounded-hier's exposure looks like there.

\begin{proposition}[Frontier sufficiency]\label{prop:sufficiency}
Fix a finite episode stream $E_0,\dots,E_{T-1}$ satisfying the frozen
stream invariants: (i) every fact's payload is observed before or at
its first required read or at its due obligation's resolution episode;
(ii) every obligation opened at episode $o$ has its gate in the
required reads of episode $\min(o+W, T-1)$, where $W$ is the
obligation window; (iii) every required read of episode $e$ is either
a current fact of $e$ or fetchable. Then the frontier policy incurs
no staleness defects, no obligation losses, and resolves every
obligation opened at $o$ with $o+W\le T-1$.
\end{proposition}

\begin{proof}
Staleness. The frontier policy never stores a believed value: a gate
enters \texttt{known} only through a fresh fact or a successful fetch,
both of which set its generation to the current episode, and eviction
only removes entries. A required read is answered from the current
facts, from a retained entry, or by a fetch (invariant (iii)); the
first two are current by construction---retained entries are updated
by every observation of their gate---and the third is current by the
reliable-retrieval assumption. Hence no read is served stale.

Obligations. An obligation opened at $o$ is inserted into the pending
set at $o$ and is exempt from eviction (the frontier retains all
pending gates). By invariant (ii) its gate is a required read at
$r=\min(o+W,T-1)$. At $r$, the gate is a current fact, or is retained
(the pending exemption keeps it in \texttt{known} from opening), or is
fetched; in all cases the resolution condition of the engine (the gate
was read validly at its due episode) is met, so the obligation
resolves. No pending obligation is ever dropped by the frontier
policy, so no loss is recorded. For $o+W\le T-1$, $r=o+W\le T-1$, so
the resolution happens within the horizon.
\end{proof}

\begin{remark}
The two directions use disjoint evidence. Necessity is measured: every
lossy-conversion policy fails on the frozen grid
(Section~\ref{sec:retire-results}). Sufficiency is proved under
invariants the frozen stream generator guarantees by construction,
which is also why the 84/84 acceptance of frontier runs is a
\emph{check} of the implementation against the proposition rather than
independent evidence for it. What the model does not prove is the
reliable-retrieval assumption itself: in a deployment where re-fetches
can fail or return stale state, the conversion rule's safety reduces
to the retrieval layer's guarantees, and bounded-hier's measured
exposure (69.0\% of due obligations and 35.9\% of footprint reads
outside its window) quantifies how much a bounded-recency design
delegates to that layer.
\end{remark}

\paragraph{Retrieval-failure sensitivity.}
The reliable-retrieval boundary admits a closed-form sensitivity
analysis. The following proposition is a declarative computation over
the already-measured operation counts, not a new experiment: the
engine's fetches never fail, so no data exists at $p>0$ and none is
simulated here. It quantifies exactly where Proposition
\ref{prop:sufficiency}'s guarantee stops, and why the conversion rule,
not a looser policy, is what survives there.

\begin{proposition}[Degradation under unreliable retrieval]\label{prop:pfail}
Fix the stream invariants of Proposition~\ref{prop:sufficiency},
replacing invariant (iii) by: each fetch succeeds independently with
probability $1-p$ for $p\in[0,1)$, and a failed fetch leaves its
required read unvalidated---recorded as a staleness defect, never
zero-imputed. Let $R(\mathcal P)$ denote the number of fetch-dependent
required reads of policy $\mathcal P$ on the stream. Then:
\begin{enumerate}[leftmargin=*,itemsep=0pt]
\item[(i)] \texttt{full-replay} has $R=0$ and incurs no
retrieval-induced defects at any $p$;
\item[(ii)] $\mathcal P$'s expected retrieval-induced defects equal
$p\,R(\mathcal P)$, and its per-run defect-free probability is
$(1-p)^{R(\mathcal P)}$;
\item[(iii)] if $R(\text{frontier})<R(\text{masking})$, frontier's
expected defects are strictly below masking's at every $p>0$, and the
ratio is $p$-independent: the measured operations gap carries over
unchanged as a defect-rate gap;
\item[(iv)] \texttt{summary-reset}'s defect count is strictly positive
at $p=0$ (109{,}761 stale reads and 385 lost obligations occurred with
retrieval fully reliable), so no retrieval layer, however perfect,
makes it safe.
\end{enumerate}
\end{proposition}
\begin{proof}
(i) \texttt{full-replay} resubmits all history and issues no fetch.
(ii) Each fetch-dependent read fails independently with probability
$p$; the count of failed reads is binomial with mean $p\,R(\mathcal P)$,
and a run is defect-free exactly when all $R(\mathcal P)$ fetches
succeed. (iii) Immediate from (ii) by linearity:
$p\,R(\text{frontier})<p\,R(\text{masking})$ for $p>0$, and the ratio
$R(\text{frontier})/R(\text{masking})$ does not depend on $p$.
(iv) These defects arise from cross-boundary information drops inside
the summary itself, not from fetches. \qed
\end{proof}

Three quantitative consequences use the long-horizon operation counts
of Table~\ref{tab:retire2-long}. First, the degradation point against
full-replay: frontier's expected excess defects reach one per accepted
run at $p\approx1/284\approx3.5\times10^{-3}$ in the clustered family
and $p\approx1/8{,}943\approx1.1\times10^{-4}$ in the independent
family; the global family has $R=0$ and is structurally immune at
every $p$. More generally, a deployment tolerating at most $\epsilon$
expected defects per run requires $p<\epsilon/R(\text{frontier})$.
Second, the ordering against masking is preserved: frontier's pooled
fetch volume is \RetwoOpsReduction{} smaller, so its expected
retrieval-induced defects are $56\%$ of masking's at every failure
rate---the conversion rule converts an operations advantage into a
correctness advantage precisely where retrieval is unreliable, and the
break-even price of Proposition~\ref{prop:breakeven} is not needed for
this direction. Third, bounded-hier enters the same formula with
$R$ equal to its out-of-window read volume, which the cohort measured
as exposure fractions (\ExpBhDueExposed{} of due obligations,
\ExpBhFootprintExposed{} of footprint reads) rather than operation
counts; its safety delegates to the retrieval layer in exactly the
sense of item (ii), with an exposure larger than frontier's
structural zero for pending obligations.

The proposition also marks the boundary of what is claimed. A real
retrieval layer fails with correlation, staleness, and byzantine
responses, not independent Bernoulli trials; retries turn $p$ into
$p^{r}$ only under independent attempts. Neither refinement is
measured here. What the model does establish is structural: the
conversion rule makes the failure surface \emph{countable}---$R$ is a
measurable property of the typed capsule---where a summary's failure
surface is silent, and item (iv) shows the difference is not a matter
of retrieval quality at all.

\begin{proposition}[Monotone break-even rule]\label{prop:breakeven}
Let $B_f,B_m>0$ and $o_f,o_m\ge 0$ be per-accepted-run regime-B means
and per-accepted-run operation means of two policies, with
$o_f<o_m$. Define $R(c)=(B_f+c\,o_f)/(B_m+c\,o_m)$ for $c\ge 0$. Then
$R$ is monotone on $c\ge0$: decreasing iff $B_f/B_m>o_f/o_m$,
increasing iff $B_f/B_m<o_f/o_m$, constant iff equal; its endpoints
are $R(0)=B_f/B_m$ and $\lim_{c\to\infty}R(c)=o_f/o_m$. For any
$\rho\in(o_f/o_m,\;B_f/B_m)$ (nonempty exactly when $R$ is
decreasing),
\begin{equation}
R(c)\le\rho\iff c\ge c_\rho,\qquad
c_\rho=\frac{\rho B_m-B_f}{o_f-\rho\,o_m}>0.
\end{equation}
If $\rho\ge B_f/B_m$ the gate holds at $c=0$; if $\rho\le o_f/o_m$ it
holds at no $c$.
\end{proposition}

\begin{proof}
$R'(c)$ has the sign of $o_f B_m-o_m B_f=B_mB_m\,(o_f/o_m-B_f/B_m)$,
which is constant in $c$, giving monotonicity with the stated
directions and endpoints. Suppose $o_f/o_m<\rho<B_f/B_m$. Solving
$B_f+c\,o_f\le\rho(B_m+c\,o_m)$ for $c$ gives
$c\,(o_f-\rho o_m)\le\rho B_m-B_f$. Here
$\rho B_m-B_f>0$ since $\rho>B_f/B_m$, and $o_f-\rho o_m<0$ since
$\rho>o_f/o_m$; dividing by the negative coefficient reverses the
inequality, so the gate is equivalent to $c\ge c_\rho$ with
$c_\rho=(\rho B_m-B_f)/(o_f-\rho o_m)>0$. If $\rho\ge B_f/B_m=R(0)$,
monotone decrease gives $R(c)\le R(0)\le\rho$ at every $c\ge0$; if
$\rho\le o_f/o_m=\lim_{c\to\infty}R(c)$, then $R(c)>\rho$ at every
finite $c$. The case $o_m=0$ forces $o_f=0$, so $R$ is constant at
$B_f/B_m$ and only the first endpoint rule applies.
\end{proof}

\begin{remark}[Deployment use]
Proposition~\ref{prop:breakeven} turns the frozen sensitivity sweep
into a decision rule. A deployment pilot that estimates the four means
$(B_f,B_m,o_f,o_m)$ and the deployment's round-trip price $c$ obtains
an exact go/no-go for the frozen 20\% promotion gate: adopt frontier
iff $c\ge c_{0.8}$, where
$c_{0.8}=(0.8B_m-B_f)/(o_f-0.8\,o_m)$ whenever
$o_f/o_m<0.8<B_f/B_m$. On the frozen cohort the pilot must actually
price its round trips: frontier's byte advantage alone ($8.0\%$
pooled) does not meet a 15\% or 20\% gate, so
$c_{0.8}\approx1{,}364$\,B pooled, $c_{0.85}\approx658$\,B pooled,
and $c_{0.8}$ is never met in the global family, where both policies
are fetch-free and the ratio is constant at $0.924$
(Section~\ref{sec:retire-exploratory}).
\end{remark}
